%%%%%%%%%%%%%%%%
%%% deut 16 %%%
%%%%%%%%%%%%%%%%

\Chapter{16}

\Verse{1}
{
	\hDtn{שָׁמוֹר אֶת חֹדֶשׁ הָאָבִיב וְעָשִׂיתָ פֶּסַח לַיהֹוָה אֱלֹהֶיךָ כִּי בְּחֹדֶשׁ הָאָבִיב הוֹצִיאֲךָ יְהֹוָה אֱלֹהֶיךָ מִמִּצְרַיִם לָיְלָה׃}
}
{
	\eDtn{
		“Observe the month of Abib, and you shall make Passover
		for YHWH, your {\God}, because in the month of Abib YHWH,
		your {\God}, brought you out from Egypt at night.
	}
}
{
%	\footnote{
%		the month of Abib. See the comment on Exod 13:4.
%	}
}

\Verse{2}
{
	\hDtn{וְזָבַחְתָּ פֶּסַח לַיהֹוָה אֱלֹהֶיךָ צֹאן וּבָקָר בַּמָּקוֹם אֲשֶׁר יִבְחַר יְהֹוָה לְשַׁכֵּן שְׁמוֹ שָׁם׃}
}
{
	\eDtn{
		And you shall make a Passover sacrifice to YHWH, your {\God},
		of the flock and herd, in the place that YHWH will choose
		to tent His name there.
	}
}
{
}

\Verse{3}
{
	\hDtn{לֹא תֹאכַל עָלָיו חָמֵץ שִׁבְעַת יָמִים תֹּאכַל עָלָיו מַצּוֹת לֶחֶם עֹנִי כִּי בְחִפָּזוֹן יָצָאתָ מֵאֶרֶץ מִצְרַיִם לְמַעַן תִּזְכֹּר אֶת יוֹם צֵאתְךָ מֵאֶרֶץ מִצְרַיִם כֹּל יְמֵי חַיֶּיךָ׃}
}
{
	\eDtn{
		You shall not eat leavened bread with it. Seven days you
		shall eat unleavened bread, the bread of degradation, with
		it, because you went out from the land of Egypt in
		haste—so that you will remember the day you went out from
		the land of Egypt all the days of your life.
	}
}
{
}

\Verse{4}
{
	\hDtn{וְלֹא יֵרָאֶה לְךָ שְׂאֹר בְּכל גְּבֻלְךָ שִׁבְעַת יָמִים וְלֹא יָלִין מִן הַבָּשָׂר אֲשֶׁר תִּזְבַּח בָּעֶרֶב בַּיּוֹם הָרִאשׁוֹן לַבֹּקֶר׃}
}
{
	\eDtn{
		And you shall not have leaven appear within all your
		borders for seven days, and none of the meat that you will
		sacrifice in the evening on the first day shall remain
		until the morning.
	}
}
{
}

\Verse{5}
{
	\hDtn{לֹא תוּכַל לִזְבֹּחַ אֶת הַפָּסַח בְּאַחַד שְׁעָרֶיךָ אֲשֶׁר יְהֹוָה אֱלֹהֶיךָ נֹתֵן לָךְ׃}
}
{
	\eDtn{
		You may not make the Passover sacrifice within one of your
		gates that YHWH, your {\God}, is giving you.
	}
}
{
}

\Verse{6}
{
	\hDtn{כִּי אִם אֶל הַמָּקוֹם אֲשֶׁר יִבְחַר יְהֹוָה אֱלֹהֶיךָ לְשַׁכֵּן שְׁמוֹ שָׁם תִּזְבַּח אֶת הַפֶּסַח בָּעָרֶב כְּבוֹא הַשֶּׁמֶשׁ מוֹעֵד צֵאתְךָ מִמִּצְרָיִם׃}
}
{
	\eDtn{
		But, rather, to the place that YHWH, your {\God}, will choose
		to tent His name: there you shall make the Passover
		sacrifice in the evening at sunset, the time when you went
		out from Egypt.
	}
}
{
}

\Verse{7}
{
	\hDtn{וּבִשַּׁלְתָּ וְאָכַלְתָּ בַּמָּקוֹם אֲשֶׁר יִבְחַר יְהֹוָה אֱלֹהֶיךָ בּוֹ וּפָנִיתָ בַבֹּקֶר וְהָלַכְתָּ לְאֹהָלֶיךָ׃}
}
{
	\eDtn{
		And you shall cook and eat it in the place that YHWH, your
		{\God}, will choose; and you shall turn in the morning and go
		to your tents.
	}
}
{
}

\Verse{8}
{
	\hDtn{שֵׁשֶׁת יָמִים תֹּאכַל מַצּוֹת וּבַיּוֹם הַשְּׁבִיעִי עֲצֶרֶת לַיהֹוָה אֱלֹהֶיךָ לֹא תַעֲשֶׂה מְלָאכָה׃}
}
{
	\eDtn{
		Six days you shall eat unleavened bread, and on the
		seventh day shall be a convocation to YHWH, your {\God}. You
		shall not do work.
	}
}
{
}

\Verse{9}
{
	\hDtn{שִׁבְעָה שָׁבֻעֹת תִּסְפּר לָךְ מֵהָחֵל חֶרְמֵשׁ בַּקָּמָה תָּחֵל לִסְפֹּר שִׁבְעָה שָׁבֻעוֹת׃}
}
{
	\eDtn{
		“You shall count seven weeks. From when the sickle begins
		to be in the standing grain you shall begin to count seven
		weeks.
	}
}
{
}

\Verse{10}
{
	\hDtn{וְעָשִׂיתָ חַג שָׁבֻעוֹת לַיהֹוָה אֱלֹהֶיךָ מִסַּת נִדְבַת יָדְךָ אֲשֶׁר תִּתֵּן כַּאֲשֶׁר יְבָרֶכְךָ יְהֹוָה אֱלֹהֶיךָ׃}
}
{
	\eDtn{
		And you shall make a Festival of Weeks for YHWH, your {\God},
		the full amount of your hand’s contribution that you can
		give insofar as YHWH, your {\God}, will bless you.
	}
}
{
}

\Verse{11}
{
	\hDtn{וְשָׂמַחְתָּ לִפְנֵי יְהֹוָה אֱלֹהֶיךָ אַתָּה וּבִנְךָ וּבִתֶּךָ וְעַבְדְּךָ וַאֲמָתֶךָ וְהַלֵּוִי אֲשֶׁר בִּשְׁעָרֶיךָ וְהַגֵּר וְהַיָּתוֹם וְהָאַלְמָנָה אֲשֶׁר בְּקִרְבֶּךָ בַּמָּקוֹם אֲשֶׁר יִבְחַר יְהֹוָה אֱלֹהֶיךָ לְשַׁכֵּן שְׁמוֹ שָׁם׃}
}
{
	\eDtn{
		And you shall rejoice in front of YHWH, your {\God}—you and
		your son and your daughter and your servant and your maid
		and the Levite who is within your gates and the alien and
		the orphan and the widow who are among you—in the place
		that YHWH, your {\God}, will choose to tent His name there.
	}
}
{
}

\Verse{12}
{
	\hDtn{וְזָכַרְתָּ כִּי עֶבֶד הָיִיתָ בְּמִצְרָיִם וְשָׁמַרְתָּ וְעָשִׂיתָ אֶת הַחֻקִּים הָאֵלֶּה׃}
}
{
	\eDtn{
		And you shall remember that you were a slave in Egypt, and
		you shall be watchful and do these laws.
	}
}
{
}

\Verse{13}
{
	\hDtn{חַג הַסֻּכֹּת תַּעֲשֶׂה לְךָ שִׁבְעַת יָמִים בְּאסְפְּךָ מִגּרְנְךָ וּמִיִּקְבֶךָ׃}
}
{
	\eDtn{
		“You shall make a Festival of Booths seven days, when you
		gather from your threshing floor and from your winepress.
	}
}
{
}

\Verse{14}
{
	\hDtn{וְשָׂמַחְתָּ בְּחַגֶּךָ אַתָּה וּבִנְךָ וּבִתֶּךָ וְעַבְדְּךָ וַאֲמָתֶךָ וְהַלֵּוִי וְהַגֵּר וְהַיָּתוֹם וְהָאַלְמָנָה אֲשֶׁר בִּשְׁעָרֶיךָ׃}
}
{
	\eDtn{
		And you shall rejoice on your festival—you and your son
		and your daughter and your servant and your maid and the
		Levite and the alien and the orphan and the widow who are
		within your gates.
	}
}
{
}

\Verse{15}
{
	\hDtn{שִׁבְעַת יָמִים תָּחֹג לַיהֹוָה אֱלֹהֶיךָ בַּמָּקוֹם אֲשֶׁר יִבְחַר יְהֹוָה כִּי יְבָרֶכְךָ יְהֹוָה אֱלֹהֶיךָ בְּכֹל תְּבוּאָתְךָ וּבְכֹל מַעֲשֵׂה יָדֶיךָ וְהָיִיתָ אַךְ שָׂמֵחַ׃}
}
{
	\eDtn{
		Seven days you shall celebrate for YHWH, your {\God}, in the
		place that YHWH will choose, because YHWH, your {\God}, will
		bless you in all your produce and all your hands’ work,
		and you shall just be happy.
	}
}
{
}

\Verse{16}
{
	\hDtn{שָׁלוֹשׁ פְּעָמִים בַּשָּׁנָה יֵרָאֶה כל זְכוּרְךָ אֶת פְּנֵי יְהֹוָה אֱלֹהֶיךָ בַּמָּקוֹם אֲשֶׁר יִבְחָר בְּחַג הַמַּצּוֹת וּבְחַג הַשָּׁבֻעוֹת וּבְחַג הַסֻּכּוֹת וְלֹא יֵרָאֶה אֶת פְּנֵי יְהֹוָה רֵיקָם׃}
}
{
	\eDtn{
		“Three times in the year every male of yours shall appear
		in front of YHWH, your {\God}, in the place that He will
		choose: on the Festival of Unleavened Bread and on the
		Festival of Weeks and on the Festival of Booths. And he
		shall not appear in front of YHWH empty-handed:
	}
}
{
}

\Verse{17}
{
	\hDtn{אִישׁ כְּמַתְּנַת יָדוֹ כְּבִרְכַּת יְהֹוָה אֱלֹהֶיךָ אֲשֶׁר נָתַן לָךְ׃}
}
{
	\eDtn{
		each according to what his hand can give, according to the
		blessing of YHWH, your {\God}, that He has given you.
	}
}
{
}

\Verse{18}
{
	\hDtn{שֹׁפְטִים וְשֹׁטְרִים תִּתֶּן לְךָ בְּכל שְׁעָרֶיךָ אֲשֶׁר יְהֹוָה אֱלֹהֶיךָ נֹתֵן לְךָ לִשְׁבָטֶיךָ וְשָׁפְטוּ אֶת הָעָם מִשְׁפַּט צֶדֶק׃}
}
{
	\eDtn{
		JUDGES “You shall put judges and officers in all your
		gates that YHWH, your {\God}, is giving you, for your tribes,
		and they shall judge the people: judgment with justice.
	}
}
{
%	\footnote{
%		judgment with justice. It should be obvious, but, regrettably, it needs
%		to be said: judgment and justice are not the same thing. Judges and
%		lawyers can be part of one of the noblest endeavors of humankind: the
%		law. When they fail to pursue justice in their performance of judgment,
%		they pervert and degrade the law, and thus they demean humankind—nothing
%		less. The same goes for clerks, bailiffs, law enforcement officers, and
%		everyone else involved in the execution of justice. They elevate or
%		lower themselves—and all of us—by the degree of their commitment to
%		justice.
%	}
}

\Verse{19}
{
	\hDtn{לֹא תַטֶּה מִשְׁפָּט לֹא תַכִּיר פָּנִים וְלֹא תִקַּח שֹׁחַד כִּי הַשֹּׁחַד יְעַוֵּר עֵינֵי חֲכָמִים וִיסַלֵּף דִּבְרֵי צַדִּיקִם׃}
}
{
	\eDtn{
		You shall not bend judgment, you shall not recognize a
		face, and you shall not take a bribe, because bribery will
		blind the eyes of the wise and undermine the words of the
		virtuous.
	}
}
{
}

\Verse{20}
{
	\hDtn{צֶדֶק צֶדֶק תִּרְדֹּף לְמַעַן תִּחְיֶה וְיָרַשְׁתָּ אֶת הָאָרֶץ אֲשֶׁר יְהֹוָה אֱלֹהֶיךָ נֹתֵן לָךְ׃}
}
{
	\eDtn{
		Justice, justice you shall pursue, so that you’ll live,
		and you’ll take possession of the land that YHWH, your
		{\God}, is giving you.
	}
}
{
%	\footnote{
%		Justice, justice. Many interpretations have been given for the
%		repetition of the word justice in this famous verse, from taking it to
%		mean that one should go to a good court (Siphre; Rashi) to midrashic and
%		Kabbalistic meanings (Ramban), as well as the most basic explanation:
%		that it is repeated for emphasis. It has also been taken to mean
%		“justice alone” (Jeffrey Tigay). And it is often said to mean that one
%		must pursue justice in a just way: one cannot use improper means to
%		achieve it. For example, police and prosecutors must not bend the law in
%		order to get a conviction. The repetition might also be taken to mean
%		that one must pursue justice for both sides: both the plaintiff and the
%		defense in a civil matter, and both the defendant and the prosecution
%		(“the people”) in a criminal matter. Most probably, though, we should
%		understand it with those who take it to be a matter of emphasis. Lacking
%		italics or exclamation points, Biblical Hebrew has few ways other than
%		repetition to express emphasis. That is the function of placing an
%		infinitive before a cognate verb (see the comment on Exod 19:12), and
%		here it is achieved by repeating the noun.
%	}
}

\Verse{21}
{
	\hDtn{לֹא תִטַּע לְךָ אֲשֵׁרָה כּל עֵץ אֵצֶל מִזְבַּח יְהֹוָה אֱלֹהֶיךָ אֲשֶׁר תַּעֲשֶׂה לָּךְ׃}
}
{
	\eDtn{
		“You shall not plant an Asherah of any wood near the altar
		of YHWH, your {\God}, that you will make.
	}
}
{
%	\footnote{
%		Asherah. The word is used for the goddess of that name or for a tree or
%		wooden pole that was associated with the worship of Asherah (or possibly
%		other deities as well).
%	}
}

\Verse{22}
{
	\hDtn{וְלֹא תָקִים לְךָ מַצֵּבָה אֲשֶׁר שָׂנֵא יְהֹוָה אֱלֹהֶיךָ׃}
}
{
	\eDtn{
		And you shall not set up a pillar, which YHWH, your {\God},
		hates.
	}
}
{
%	\footnote{
%		you shall not set up a pillar. Moses set up pillars himself (Exod 24:4).
%		So did Jacob (Gen 28:18; 32:45; 35:14,20). So why does Moses forbid it
%		now? Rashi and some modern commentators have thought that pillars were
%		originally acceptable but came to be associated with pagan worship and
%		so were abandoned. That may or may not be the actual historical
%		development behind these changes, but that does not seem to be what is
%		presented here in the text, because Moses moves from making them to
%		saying that God hates them, all within the forty years in the
%		wilderness. Perhaps we should understand it to mean that Moses is
%		prohibiting the people from making pillars after the formal altar and
%		Temple are built in the land, because then an altar might be taken as an
%		alternative locus of God. Here in Deuteronomy Moses repeatedly
%		emphasizes that there can be only one place of worship, one structure in
%		which God’s name is housed. Thus, even Joshua sets up a large stone,
%		which seems like a pillar (although the word is not used there; Josh
%		24:26), but this may be before the establishment of an official central
%		place of worship.
%	}
}
